\documentclass[a4paper,10pt]{article}
\usepackage[utf8]{inputenc}
\usepackage{geometry}
\usepackage{enumitem}
\usepackage[hidelinks]{hyperref}
\usepackage{titlesec}
\usepackage{changepage}

\geometry{margin=0.65in}
\setlist[itemize]{left=0pt, itemsep=0pt, parsep=0pt, topsep=4pt}
\titleformat{\section}{\normalfont\Large\bfseries\centering}{}{0pt}{}
\titleformat{\subsubsection}{\normalfont}{}{0pt}{}
\titlespacing*{\subsection}{0pt}{10pt}{5pt}
\titlespacing*{\subsubsection}{0pt}{7pt}{4pt}

\begin{document}
\pagestyle{empty}

\begin{center}
  \section*{\LARGE Robert Martinez Marenco}
  \href{mailto:rdmm404@gmail.com}{rdmm404@gmail.com} |
  \href{https://github.com/rdmm404}{github.com/rdmm404} |
  \href{https://linkedin.com/in/rdmm404}{linkedin.com/in/rdmm404}

  Permanent Resident in Toronto, ON
\end{center}

\hrule
\vspace{4pt}

% Professional Summary
\subsection*{Summary}
Senior Software Engineer with 4+ years of experience building scalable cloud microservices and full-stack web applications.
Skilled at translating business requirements into reliable product workflows using Python, Go, TypeScript, React,
GraphQL, and modern CI/CD practices.
% Work Experience
\subsection*{Work Experience}

\subsubsection*{%
  \textbf{\href{https://www.achievers.com/}{Achievers}, Toronto}
  \hfill\textit{November 2022 - Present}
}
\begin{adjustwidth}{1em}{0pt}
\noindent\textbf{Senior Software Engineer}
\hfill\textit{October 2026 - Present}\\
\textbf{Intermediate Software Engineer}
\hfill\textit{October 2025 - October 2026}
\begin{itemize}
  \item Develop and maintain Python and Go microservices for a global B2B SaaS platform using GCP, gRPC, Pub/Sub, and Kubernetes (GKE).
  \item Co-designed and implemented a resilient top-up checkout workflow with Python backend flows, Go-based Avalara tax integration, Stripe payments, and saga-style distributed transaction coordination.
  % Add captured revenue metric once available.
  % Interview details: fallback polling jobs, checkout queue state machine, Pub/Sub redelivery, DLT management,
  % UUID deduplication, simultaneous checkout locks across devices, abort/refund paths.
  % TODO: Consider splitting into design + implementation bullets for payment/platform-focused roles.
  \item Helped simplify the Elasticsearch indexing pipeline through architecture evaluation, Go batch workers, queue tuning, and Airflow/Dataflow changes; reduced cluster CPU by approximately 30\% and memory pressure by approximately 20\%, with search latency under 200ms.
  % \item Partnered with SRE to optimize Elasticsearch index mappings by removing unused indexed fields, correcting data types,
  %   and preventing unwanted extra fields, reducing UAT cluster memory pressure by 50\%+ ahead of production rollout.
  \item Built a RAG pipeline and MCP server for vendor API documentation retrieval and JSON config validation, cutting integration timelines from 4 weeks to 3 and helping the team exceed its annual integration goal by 20\%.
  \item Improved New Relic observability for Elasticsearch data sync queues with custom queue-state events, backlog alerts, and dashboards; the implementation pattern was later standardized in the shared batch processing library.
  \item Mentored 2 co-op engineers through onboarding and project delivery with code reviews, pair programming, and backend technical guidance.
  % Optional AI leadership bullet:
  % \item Contributed marketplace-focused AI skills and checkout triage runbooks to an internal skill marketplace,
  %   and hosted team knowledge-transfer sessions to improve AI-assisted development and production support workflows.
  % TODO: Uncomment once shipped.
  % \item Improved checkout observability by tagging New Relic transactions with order metadata and terminal outcomes,
  %   enabling alerts for abnormal aborted-order rates and dashboards for checkout volume, success, and failure trends.
  % Future/search architecture bullet: Sphinx deprecation project, search architecture tradeoffs,
  % Elasticsearch migration decisions, and implementation impact once shipped/measured.
\end{itemize}

\smallskip
\noindent\textbf{Software Engineer}
\hfill\textit{November 2022 - October 2025}
\begin{itemize}
  \item Developed React components across Single SPA micro frontends, collaborating with Product and UX on Figma-based workflows while integrating GraphQL gateway endpoints with backend microservices.
  \item Maintained highly available order processing and vendor integration services for a marketplace platform handling 10+ orders/min (\$10M+ 2024 revenue), resolving critical production issues within SLAs.
  \item Integrated external APIs including NetSuite and Akeneo to streamline product merchandising workflows, improving pipeline efficiency by 66.6\%; optimized MySQL and Redis-backed service performance.
\end{itemize}
\end{adjustwidth}

\subsubsection*{%
  \textbf{Web Developer} at \href{https://www.tecnoglass.com/en/}{\textbf{Tecnoglass}}, Colombia
  \hfill
  \textit{November 2021 - November 2022}
}
\begin{itemize}
  \item Developed full-stack web applications using Python, Django, JavaScript, and PostgreSQL to reduce paperwork and optimize business processes; built a centralized Flask REST API for LDAP-backed authentication.
  \item Designed and implemented a LoRaWAN temperature and humidity monitoring system for glass assembly plants, including hardware selection, installation, and a real-time Django Channels dashboard.
\end{itemize}

% Education
\subsection*{Education}
\subsubsection*{\textbf{B.Sc. in Electronic Engineering, cum laude} \hfill \textit{March 2023}}
\noindent Universidad del Norte, Colombia

% Skills
\subsection*{Skills}
\begin{itemize}
  \item \textbf{Languages:} Python, Go, JavaScript, SQL, TypeScript, PHP, HTML, CSS.
  \item \textbf{Frameworks:} React, GraphQL, Flask, FastAPI, Django.
  \item \textbf{Cloud \& DevOps:} Google Cloud Platform (GCP), AWS, DroneCI, GitHub Actions, Airflow.
  \item \textbf{Technologies:} MySQL, PostgreSQL, MongoDB, Redis, Elasticsearch, gRPC, Pub/Sub, Docker, Kubernetes, New Relic.
\end{itemize}

% TODO: Projects section
\end{document}
